% ===================================================================================================================
% Section IV body (paste under \section{TIDES Architecture} \label{sec:arch}).
% Every number is the RTL (hw/rtl/ours/ours_core.v) / cycle model (ours.py Cfg) configuration of Table I.
% TODO architecture overview figure (Figures/TIDES_ppu): \ref{fig:arch}.
% ===================================================================================================================
TIDES executes the spiking GEMM of Eq.~\eqref{eq:gemm} layer by layer on a weight-stationary systolic array of
$R \times C$ PEs with $T$ adders each. It consists of a spike buffer and a weight buffer, a pre-processing unit
that applies TIP, a group buffer, the PE array, a post-processing unit at the bottom of the array, an output
buffer and an array of LIF units. Only the pre- and post-processing units are specific to DPTP; the PEs perform
label-indexed accumulation and nothing else. The evaluated configuration uses $R = C = 4$, $T = 8$,
$G_{\max} = F_{\max} = Q = 4$ (Section~\ref{sec:eval-method}).

\subsection{Execution Flow}
\label{sec:arch-flow}
A layer is split into output tiles of 64 filters and 128 output positions, whose partial sums stay in the output
buffer until the tile is complete, and each output tile into input tiles $\mathcal{K}$ of $R$ inputs. For every
input tile, the packed weight columns of Algorithm~\ref{alg:filter} are loaded into the PEs, the pre-processing
unit packs the tile's spike trains with Algorithm~\ref{alg:tip}, and every group streams through the array once
per pass over the packed columns, $C$ columns at a time. The post-processing unit adds each finished partial sum
into the output buffer. After the last input tile, the LIF units turn the accumulated currents into the output
spike trains of the tile, which are written back to the spike buffer as the input of the next layer. Weight
loading, packing and streaming are pipelined, so in steady state the array receives one group per cycle.

\subsection{Pre-processing Unit}
\label{sec:arch-pre}
The pre-processing unit implements Algorithm~\ref{alg:tip} in hardware. It reads, for each output position of an
output tile, its $R$ spike trains on the current input tile from the spike buffer and forms the $RT$-bit mask
$\mathit{mask}_m$; positions with an all-zero mask are dropped. A table of $Q$ open groups holds, per group, its
accumulated mask, the IDs of its members and their count. Each incoming mask is ANDed with the $Q$ group masks in
parallel, and a first-fit priority encoder selects the oldest group with a zero result; the mask is ORed into
that group and the position's ID appended. A group that reaches $G_{\max}$ members, or the oldest group when no
group fits and all $Q$ are open, is emitted to the group buffer. The unit has eight independent lanes, each
packing a different input tile, so packing runs ahead of the array.

An emitted group is one 128-bit group-buffer entry: its $RT = 32$ spike bits, a 2-bit member label per bit, the
7-bit IDs of its up to four members, and a count. The group buffer is double-buffered (2 $\times$ 2\,KB), so the
unit fills one half while the array drains the other.

\subsection{PE Array}
\label{sec:arch-array}
Packed trains move right along the rows and partial sums move down the columns. Row $k$ of a group enters the
array one cycle after row $k-1$, so a group's partial sums meet its packed train in every row. Each row word
carries a valid bit, a pass parity bit, the $T$ spike bits of $\hat{\mathbf{s}}_{G,k}$ and their $T$ member labels
$\ell_{G,k}[t]$; each column carries the bundle of $F_{\max} G_{\max}$ partial sums of the group, $T$ values of
10 bits each, together with one touch bit per partial sum. Ten bits suffice because a partial sum in the array
accumulates at most $R$ INT8 products before it leaves.

\textbf{Processing element.} A PE holds one weight in an active register: an 8-bit value, its 2-bit filter label
and a pass parity bit. A shadow register receives the weight of the next pass while the current pass streams, and
replaces the active register when a row word with the other parity arrives, so loading new weights does not stall
the array. The PE has one lane per timestep. Lane $t$ concatenates the filter label with $\ell_{G,k}[t]$ into a
4-bit slot index, reads that partial sum from the incoming bundle, adds the weight, and writes the result back to
the same slot, setting its touch bit. The write is enabled only when the row word is valid, $\hat s_{G,k}[t] = 1$
and the weight is non-zero; all other slots pass to the PE below unchanged, and the row word passes to the PE on
the right. Because a packed train carries at most one spike per timestep and a PE at most one weight, the $T$
lanes write to distinct partial sums and need no arbitration. The PE contains no comparator, prefix-sum circuit or
search logic: its datapath is $T$ adders, the slot multiplexers and the pipeline registers.

\subsection{Post-processing Unit}
\label{sec:arch-post}
Each column ends in one lane of the post-processing unit. When a group leaves the last row, only its touched
partial sums produce entries. A lane restores each entry's address from its labels: the filter label indexes the
column's filter list and the member label indexes the group's member IDs, both carried in a sideband that is
delayed to arrive with the bundle. Each entry, the 6-bit filter ID, the 7-bit position ID and the $T$ partial
sums, enters an 80-entry FIFO. The FIFO is split into 16 banks, so the up to 16 entries of one group are written in
the same cycle. An arbiter retires up to four entries per column per cycle with a two-stage read-modify-write: it
reads the 24-bit sums of all $T$ timesteps of the entry's (filter, position) from the output buffer, adds the
partial sums and writes them back, and keeps an address in flight out of the next issue slot to avoid a
read-after-write hazard. The FIFOs absorb the burst of entries of each group, and the controller issues a group
only if every column can accept its entries, so no entry is lost.

\subsection{Output Buffer and LIF Units}
\label{sec:arch-lif}
The 192\,KB output buffer holds one entry of $T$ 24-bit currents per (filter, position) of the output tile. After
the last input tile, 16 LIF units read the entries that received at least one partial sum and evaluate
Eq.~\eqref{eq:lif} with $\lambda = 1/2$. Each unit is pipelined over the timesteps, two timesteps per stage, and
accepts one entry per cycle, producing a $T$-bit output spike train. Entries that received no partial sum carry no
current, emit no spike and are skipped.

\subsection{Why Weight-Stationary}
\label{sec:arch-ws}
TIP forms groups per input tile, so a group needs to be conflict-free only on the $R$ inputs of that tile, and
the next tile regroups the positions. A weight-stationary array realizes exactly this: weights change once per
input tile, and partial sums leave the array after $R$ rows. In an output-stationary array, a PE keeps its outputs
for all $K$ inputs of the layer, so a group would have to be conflict-free on every input, and two positions almost
always collide somewhere among hundreds of inputs. The cost of weight stationarity is that partial sums of
consecutive input tiles meet in the output buffer, which the post-processing unit accumulates; Section~\ref{sec:eval}
compares the three mappings.
